
\section{Disentangling Morphing from Shortcuts}
\label{sec:controlled_eval}

Conventional single-image morphing attack detection (S-MAD)
evaluates whether a detector can distinguish morphed faces
from genuine photographs. However, these two classes may
differ in more than the presence of multiple identities.
Morph generation can introduce geometric transformations,
blending artifacts, generative reconstruction traces, and
differences in image processing or storage. Consequently,
high classification accuracy does not reveal which of these
properties drives the detector's predictions.

We introduce a framework of \emph{matched controls} to
investigate these alternative explanations. Rather than
treating morph detection as a single binary comparison,
we construct pairs of image populations that share
increasingly similar generation procedures while differing
in selected aspects of their construction.

\subsection{Morphing and Nuisance Factors}
\label{sec:morph_factors}

Let $x_i$ and $x_j$ denote photographs of individuals
$i$ and $j$, and let
\begin{equation}
    x_{ij}^{m} = M_e(x_i,x_j), \qquad i \neq j,
\end{equation}
denote a morph produced by engine $e$.
A conventional evaluation compares $x_{ij}^{m}$ with
genuine photographs $x_i$.

This comparison simultaneously changes several properties
of the image. Landmark-based morphing methods introduce
geometric warping and blending, while GAN- and diffusion-based
approaches additionally involve generative reconstruction
\cite{zhang2021mipgan,damer2023mordiff}.
Moreover, morphs and genuine images may differ in their
resizing, compression, alignment, and storage histories.

Let $f(x)$ be a detector's morph score and $P_h(x)$
an image subjected to processing history $h$.
Our objective is to determine which image properties
are sufficient to produce a high $f(x)$, and whether
discrimination persists when the compared populations
undergo similar transformations.

\subsection{A Hierarchy of Matched Controls}
\label{sec:matched_controls}

We consider five complementary comparisons, progressing
from processing artifacts to differences in contributor
identity. Each addresses a specific alternative explanation
for successful morph detection.




\paragraph{Processing controls.}
We compare genuine photographs $x_i$ with processed
versions $P_h(x_i)$, including images subjected to
downsampling, enlargement, or changes in storage history.
Since identity content is unchanged, a substantial
difference in detector responses reveals sensitivity to
processing alone. Such sensitivity does not, however,
establish that processing explains every correct morph
prediction.

\paragraph{Generative reconstruction controls.}
For GAN- and diffusion-based morphing, we compare camera
photographs with images reconstructed from a single
individual using the corresponding generative pipeline.
These controls contain no second identity but retain
generation-related artifacts. A high morph score on
single-person reconstructions indicates that the rendering
process itself can trigger detection.

\paragraph{Warp-only controls.}
For landmark-based morphing, we construct images using
the geometric transformations of the morphing procedure
without blending two photographs. Comparing these
warp-only controls against blended morphs tests whether
blending introduces a detectable signal beyond geometric
warping. However, differences in noise, local variance,
and interpolation may remain, so successful discrimination
does not necessarily indicate sensitivity to identity mixing.

\paragraph{Same-person morph controls.}
We construct morphs from two different photographs of the
same individual:
\begin{equation}
    x_{ii'}^{m} = M_e(x_i,x_i'),
\end{equation}
where $x_i$ and $x_i'$ depict the same identity.
These images undergo the morphing procedure but do not
combine two distinct identities. Comparing
$x_{ij}^{m}$ with $x_{ii'}^{m}$ therefore tests whether
detector scores distinguish two-person manipulations
from similar same-person manipulations.

Importantly, same-person and two-person morphs may still
differ in expression, pose, texture, and blending strength.
This comparison measures contributor-dependent
discrimination, rather than proving that the detector
recognizes two identities.

\paragraph{Face-splice controls.}
Whole-frame morphs may contain artifacts outside the face,
including changes to hair and background. To investigate
local manipulation, we construct face-spliced morphs by
inserting the manipulated facial region into a geometrically
matched host image.

We consider three constructions:
(i)~two-person face splices, containing facial information
from different identities;
(ii)~same-person face splices, constructed from two images
of one individual; and
(iii)~seam-only controls, which retain an analogous
insertion procedure without introducing a second identity.


\paragraph{Splice construction.}
We construct controls using the same 106-point landmark
geometry and averaged target shape as the corresponding
morph. The insertion region $H$ is the convex hull of
the averaged facial landmarks, rounded to integer
coordinates. A two-person splice copies pixels from
the two-person morph into its corresponding warp-only
image, while a same-person splice uses the same-person
morph and its warp-only counterpart. Seam-only controls
instead insert same-person warp-only pixels into the
two-person warp-only image. All constructions preserve
the image outside $H$ and use direct pixel copying,
without additional interpolation. Thus, insertion
geometry is shared, although local texture and blending
statistics may still differ.


\subsection{Matched Construction and Evaluation}
\label{sec:control_protocol}

To ensure consistent comparisons, we reconstruct the
landmark-based morphing procedures used in MAD22 and
FRLL-Morphs, including their warping geometry, alignment,
and post-processing. The corresponding controls are
generated with consistent construction procedures.
Further implementation and provenance details are
provided in the supplementary material.

We evaluate detector responses using complementary
metrics. Area under the ROC curve (AUC) measures the
ranking of matched image populations, while equal error
rate (EER) enables comparison with conventional MAD
benchmarks. We additionally report false-alarm and
attack-miss rates at fixed decision thresholds, since
ranking performance alone does not establish reliable
operational behavior.

\paragraph{Controlling manipulation strength.}
Two-person and same-person morphs can differ in the
magnitude of image changes introduced by blending.
To reduce this imbalance, we employ a splice-derived
change proxy. For a splice $S$, its corresponding
outside-image construction $C$, and facial insertion
region $H$, define
\begin{equation}
    D(S,C) =
    \frac{1}{3|H|}
    \sum_{p \in H}
    \sum_{c=1}^{3}
    |S_{pc}-C_{pc}|.
    \label{eq:splice_difference}
\end{equation}

Here, $D$ measures the mean absolute RGB change within
the insertion region. We stratify whole-frame comparisons
into quintiles of $\log D$, retaining strata with at
least ten examples from each class and requiring retained
strata to cover at least half of the evaluated examples.


The proxy is computed on stored $512$-pixel RGB images
with intensities in $[0,255]$. For whole-frame
comparisons, each image is assigned the $D$ value of
its corresponding splice construction; this adjusts
for splice-derived change rather than directly matching
whole-frame blending strength. Within each retained stratum $q$, we compute the
discrimination AUC $a_q$ and aggregate it as
\begin{equation}
    \mathrm{AUC}_{D} =
    \frac{\sum_q n_q^{+}n_q^{-}a_q}
         {\sum_q n_q^{+}n_q^{-}},
    \label{eq:stratified_auc}
\end{equation}
where $n_q^{+}$ and $n_q^{-}$ are the numbers of
two-person and same-person examples in stratum $q$.
The resulting AUC is aggregated across retained strata,
weighted by the numbers of positive--negative pairs.
This reduces differences in measured manipulation
strength but does not match all low-level image properties.

% \paragraph{Statistical interpretation.}
% Whenever possible, comparisons use the same source
% photographs or contributor pairs. Uncertainty estimates
% account for dependencies among images sharing identities,
% hosts, or morphing partners through component-level
% bootstrap resampling.

% We also evaluate simple image-statistic references to
% determine whether controlled discrimination could be
% explained by low-level properties. Such references are
% interpreted according to their training and evaluation
% populations and are not automatically treated as
% out-of-domain competitors to pretrained detectors.

% Crucially, passing an individual control does not prove
% that a detector relies on identity-specific evidence.
% Each comparison rules out or reduces only particular
% alternative explanations.

\paragraph{From evaluation to training.}
The controls introduced above also provide counterexamples
for detector training. If processing history, geometric
warping, or insertion artifacts occur under both class
labels, these cues become less reliable predictors of
morphing. In the following section, we use this principle
to construct a matched-control training protocol that
aims to reduce shortcut dependence while improving
discrimination between two-person and same-person
manipulations.







%%%%%%%%%%%%%%%%%%%%%%%%%%%%%%%%%%%%%%%%%%%%%%%%%%%%%%%%%%%%%%%%%%%%%%%%%%%%%%%%%%%%%%%%%%%%%%%%%%%%%


\section{Shortcut-Aware Training with Matched Controls}
\label{sec:matched_training}

The controls introduced in Section~\ref{sec:controlled_eval}
serve not only to identify shortcut dependencies but also
to construct training examples that challenge them.
Conventional morph detection training can associate image
processing, geometric warping, or blending artifacts almost
exclusively with the morph class. As a result, these cues
may become predictive even when they provide no evidence
that two distinct identities contributed to the image.

We investigate a simple intervention: introduce matched
examples that share such artifacts across both class labels.
Our goal is not to enforce complete invariance to image
processing, but to reduce the sufficiency of particular
shortcuts while encouraging discrimination between
two-person morphs and comparable single-person
manipulations. Unlike approaches that introduce new
architectures or specialized objectives, we modify only
the training data and keep the detector and optimization
procedure fixed.

\subsection{Training Data Construction}
\label{sec:training_data}

We construct four principal training configurations,
progressively incorporating the controls from
Section~\ref{sec:matched_controls}.

\paragraph{Plain baseline (A).}
The baseline consists of genuine photographs labeled
as bona fide and conventional whole-frame morphs
constructed from two different identities labeled
as attacks. This represents standard supervised
morph detection training.

\paragraph{Processing-matched controls (B).}
We augment A with three types of negative examples:
(i) geometrically warped photographs without blending,
(ii) individual warped input halves extracted from the
morph construction, and
(iii) copies of genuine images stored at reduced resolution.

Importantly, reduced-resolution copies also occur in
the positive class through corresponding two-person
morphs. Consequently, storage history is no longer
exclusively associated with the attack label.
Warp-only examples and individual halves similarly
expose geometric and preprocessing artifacts without
requiring a second identity.

\paragraph{Face-spliced positives (Sp).}
Whole-frame morphs may alter regions outside the face,
allowing detectors to exploit changes in hair, background,
or image boundaries. To encourage sensitivity to local
facial manipulation, we extend B with two-person
face-spliced morphs.
Each example inserts the morphed facial region into a
geometrically matched host image, preserving the
surrounding image content. These examples are labeled
as attacks because they combine information from
two distinct identities.
However, splice positives alone may introduce a new
shortcut: the network can learn to associate the
presence of a splice boundary with the positive label.

\paragraph{Full matched-control protocol (Sb).}
To address this ambiguity, we supplement Sp with
two complementary negative constructions:
(i) same-person face splices, created from two
photographs of the same individual, and
(ii) seam-only controls, which reproduce an
insertion procedure without introducing facial
content from a second identity.

These negatives retain aspects of the manipulation
process while removing the intended two-person
identity change. Consequently, merely detecting
an insertion seam or a local manipulation is
insufficient to reproduce the training labels.
We treat these examples as negative controls for
the specific objective of detecting two-person
morphs, rather than as unmodified photographs.
This distinction is important: a same-person splice
is manipulated, but it does not constitute a
two-identity morph.

\subsection{Controlled Training Comparisons}
\label{sec:training_comparisons}

The progression A--B--Sp--Sb allows us to examine
the effects of processing controls, face-spliced
positives, and matched splice negatives.
However, simply removing a training component may
also change the number of examples, class balance,
and effective exposure during optimization.
We therefore introduce count-balanced comparisons
that hold the number of training rows and positive
examples fixed within each pair.

\paragraph{Matched negatives with splice positives.}
Our primary comparison evaluates Sb against SpV,
a variant of Sp in which repeated ordinary negative
examples replace the matched splice negatives.
Both configurations contain 8,386 balanced training
rows and the same positive examples, including
their repetitions.

The difference is that Sb introduces 312 additional
unique matched-negative images, whereas SpV repeats
existing negatives. Thus, the comparison evaluates
the combined effect of matched-negative content and
additional image diversity, rather than isolating
the semantic effect of the controls alone.

\paragraph{Matched negatives without splice positives.}
To examine whether matched negatives are beneficial
without two-person face-spliced examples, we additionally
compare BN with BV. Both are derived from B and contain
8,386 balanced training rows, but BN incorporates
matched splice negatives whereas BV repeats ordinary
negative examples.

Together, these comparisons examine whether matched
negatives provide useful supervision in the presence
or absence of corresponding two-person splice positives.
They do not, by themselves, isolate a causal interaction
between the two components, since the positive training
distributions differ across pairs.

\paragraph{Component ablations.}
We further examine the individual contributions of
warp-only negatives, warped input halves, and
reduced-resolution copies by removing each component
from the full protocol. Additional variants test the
effect of adding stored-small copies without the
remaining controls. Their construction and training
counts are detailed in the supplementary material.
